\documentclass[UTF8]{ctexart}
\makeatletter
\def\input@path{{../../}{../}}
\makeatother
\usepackage{researchnotes}
\title{江苏高考数学附加题压轴题汇编}
\author{}
\date{检索范围：2008—2019 年}
\begin{document}
\maketitle
\tableofcontents

\section{编写范围与阅读方法}\label{sec:intro}
\subsection{两个不同的起点}
江苏高考自主命题始于 2004 年，而本书研究的“160 分数学主卷加 40 分数学附加题”始于 2008 年，二者不能混为一谈。2008 年考试说明明确规定：数学必做部分满分 160 分、考试时间 120 分钟，选测物理的考生另做满分 40 分、考试时间 30 分钟的附加题。因此，按截至 2019 年的时间上限，本书高考部分的对象是 2008—2019 年共 12 年的数学附加题末题。2004—2007 年不具备这一试卷结构，不将其主卷末题混入；2020 年亦不在本次时间上限内。制度依据见文末来源~\ref{src:history}、\ref{src:scheme}。

这里的“压轴题”按卷面位置界定，即附加题最后一道必做解答题，并不预先断言它一定是全卷最难的题。原高考卷通常将它编号为第 23 题；一些题库把第 21 题的各个选做小题拆开重新编号，末题可能变成第 25 或第 26 题。辨认时以题干及其处于附加题末尾的位置为准，不只搜索题号。

\subsection{市级卷的收录口径与核验界限}
市级部分以南京、无锡、徐州、常州、苏州、南通、连云港、淮安、盐城、扬州、镇江、泰州、宿迁 13 个设区市为检索对象，检索同一时期的一模、二模、三模，以及具有相应功能的市级调研测试。市级联考只收录一道题，通过城市索引指向同一题解。原题只写“期末检测”“教学情况调研”时，保留原名；未经原卷或同期通知确认，不擅自把它改叫一模。单校测试、区县测试、商业机构预测卷和仅挂有“江苏”名称的练习卷均不自动算作市级模拟卷。

本书是带来源登记的汇编稿，\keyterm{并非已经收全的 13 市历年全套题库}。正文题目与文末检索登记分开：正文给出能够确认数学条件的题意重述、独立推导和方法说明；只找到目录、残缺公式、失效下载地址或存在归属冲突的材料进入待核清单，不能计入已完成题解数。“未核得”表示本次检索未获得足够证据，不表示该市当年没有举办该轮考试。所谓 $12\times13\times3=468$ 只是检索槽位上限，不是已经证明存在的 468 份独立试卷。

\subsection{先识别结构，再实施计算}
这组题的共同困难往往是把题目转化为可计算对象。计数题应先回答“一个对象由哪些选择唯一确定”，概率题应先写出等可能样本空间，递推题应解释从规模 $n-1$ 到 $n$ 增加或删除了什么，高阶导数题则应寻找低阶关系在求导后是否保持形式。每题的“思路”用于解释这一选择；“解答”给出可核查步骤；“复盘”指出方法的适用条件与容易遗漏的边界。除标为补充的内容外，不用另造练习充当市级原题。

\section{常用工具}\label{sec:tools}
以下工具只承担后文解题所需的准备作用。组合数采用 $\binom nr$，对整数 $n\geq0$，规定 $r<0$ 或 $r>n$ 时为零。$\lfloor x\rfloor$、$\lceil x\rceil$ 分别表示向下、向上取整。

\begin{lemma}[组合数的相邻关系与逐项累加]\label{lem:binomial}
对整数 $n\geq1$、$0\leq r\leq n$，有
\[
\binom nr=\binom{n-1}r+\binom{n-1}{r-1},\qquad
(r+1)\binom{n+1}{r+1}=(n+1)\binom nr,
\]
第二个等式在 $n=0,r=0$ 时也成立。此外，对整数 $N\geq r\geq0$，有
\begin{equation}\label{eq:hockey}
\sum_{j=r}^{N}\binom jr=\binom{N+1}{r+1}\qquad(N\geq r).
\end{equation}
\end{lemma}
\begin{proof}
第一个等式按某个指定元素是否被选中分成互斥的两类；第二个等式可由阶乘定义直接约去公共因子。再由第一个等式得 $\binom jr=\binom{j+1}{r+1}-\binom j{r+1}$，从 $j=r$ 累加至 $N$，相邻项抵消，得到式~\eqref{eq:hockey}。
\end{proof}

\begin{lemma}[插入最大元素与逆序数]\label{lem:insert}
一个 $n$ 元排列的逆序数（inversion number）是满足 $i<j$ 且 $a_i>a_j$ 的下标对数。把最大元素 $n$ 插入一个 $n-1$ 元排列，若插入后它的右侧有 $j$ 个元素，则新增逆序数恰为 $j$，旧元素之间的逆序数不变。
\end{lemma}
\begin{proof}
旧元素的相对顺序没有变化。最大元素与其左侧元素不构成逆序，而与右侧每一个元素恰构成一个逆序。删除最大元素又能唯一恢复原排列及插入位置，所以这一构造没有重复也没有遗漏。
\end{proof}

\begin{lemma}[相同成功率的独立试验]\label{lem:binomial-law}
若 $N$ 次试验相互独立，每次成功概率均为 $p$，成功次数 $Y$ 服从二项分布（binomial distribution），即
\[
\mathbb P(Y=j)=\binom Njp^j(1-p)^{N-j},\quad 0\leq j\leq N,
\qquad \mathbb E Y=Np.
\]
\end{lemma}
\begin{proof}
恰成功 $j$ 次需要选择 $j$ 个成功位置，共 $\binom Nj$ 种；每一种位置配置的概率由独立性等于 $p^j(1-p)^{N-j}$。若以 $I_i$ 表示第 $i$ 次成功的示性随机变量，则 $Y=\sum I_i$，故 $\mathbb E Y=\sum\mathbb E I_i=Np$。这里期望的可加性本身不要求独立；独立性用于前面的分布公式。
\end{proof}

\section{2008—2019 年高考附加题末题}\label{sec:gaokao}
\subsection{2008 年：二项式中的权与微积分}\label{sec:g2008}
\textbf{来源与校勘。}见来源~\ref{src:g2008}，原卷末题为第 23 题，所用转载将它重排为第 26 题。题目前以余弦倍角公式求导得到正弦倍角公式，提示对恒等式施加运算。转载的最后一个求和式存在下限或右端排印疑点；下面以从 $k=0$ 开始的正确完整形式讲解，并同时给出从 $k=1$ 开始时的数值，不将未经核清的转录式原样作为命题。

\begin{example}[题意与经校验的恒等式]
利用 $(1+x)^n=\sum_{k=0}^{n}\binom nkx^k$，对整数 $n\geq2$ 证明
\[
n\bigl((1+x)^{n-1}-1\bigr)=\sum_{k=2}^{n}k\binom nkx^{k-1}.
\]
对整数 $n\geq3$，进一步证明
\[
\sum_{k=1}^{n}(-1)^kk\binom nk=0,\qquad
\sum_{k=1}^{n}(-1)^kk^2\binom nk=0,
\qquad
\sum_{k=0}^{n}\frac{1}{k+1}\binom nk=\frac{2^{n+1}-1}{n+1}.
\]
\end{example}
\textbf{思路。}求和中出现 $k$，与单项式 $x^k$ 求导时产生的因子相吻合；出现 $k^2$ 时，将它拆成 $k(k-1)+k$；出现 $1/(k+1)$ 时，想到 $x^k$ 的积分。符号 $(-1)^k$ 提示在导数恒等式中代入 $x=-1$。
\begin{proof}
二项式展开式两端求导得
\[
n(1+x)^{n-1}=\sum_{k=1}^{n}k\binom nkx^{k-1}.
\]
把 $k=1$ 的项 $n$ 移到左边即为第一问。代入 $x=-1$，并把等式两边乘以 $-1$，得第一条交错和。再次求导，有
\[
n(n-1)(1+x)^{n-2}=\sum_{k=2}^{n}k(k-1)\binom nkx^{k-2}.
\]
因为 $n\geq3$，在 $x=-1$ 时左端为零；右端的 $(-1)^{k-2}$ 等于 $(-1)^k$。与第一条交错和相加，就得到 $k^2$ 加权的交错和。

最后，把有限和逐项在 $[0,1]$ 上积分，得
\[
\sum_{k=0}^{n}\binom nk\int_0^1x^k\,\mathrm dx
=\int_0^1(1+x)^n\,\mathrm dx
=\left.\frac{(1+x)^{n+1}}{n+1}\right|_0^1
=\frac{2^{n+1}-1}{n+1}.
\]
若下限改为 $1$，应减去 $k=0$ 的项 $1$，所以正确右端为 $(2^{n+1}-n-2)/(n+1)$。这一校验不能省略。
\end{proof}
\textbf{复盘。}也可利用 $\binom nk/(k+1)=\binom{n+1}{k+1}/(n+1)$，将积分恒等式直接化为二项式系数和。求导中 $n\geq3$ 的条件保证二阶导数在 $x=-1$ 时为零，不能把它机械推广到 $n=2$。

\subsection{2009 年：用判别式划分整数点}\label{sec:g2009}
\textbf{来源。}见来源~\ref{src:g2009}。
\begin{example}[题意重述]
从 $1,2,\ldots,n$ 中分别等可能、独立地选取 $a,b$，允许 $a=b$。记 $T_n$ 为使方程 $x^2+2ax+b=0$ 有实根的有序数对 $(a,b)$ 的个数，$P_n$ 为相应概率。求 $T_{n^2}$、$P_{n^2}$，并证明对整数 $n\geq2$ 有 $P_n>1-1/\sqrt n$。
\end{example}
\textbf{思路。}实根条件就是 $b\leq a^2$。固定 $a$ 后，允许的 $b$ 的个数为 $\min(n,a^2)$，所以只需在 $a$ 穿过 $\sqrt n$ 处断开求和。第二问不必算出所有取整情形，只估计无实根的数对即可。
\begin{proof}
共有 $n^2$ 个等可能的有序数对，故
\begin{equation}\label{eq:g2009-count}
T_n=\sum_{a=1}^{n}\min(n,a^2),\qquad P_n=T_n/n^2.
\end{equation}
为避免把题目下标中的平方与样本空间大小混淆，第一问先令选数上界为 $N=n^2$。当 $1\leq a\leq n-1$ 时有 $a^2$ 种选法；当 $n\leq a\leq n^2$ 时有 $n^2$ 种选法。因此
\[
T_{n^2}=\frac{n(n-1)(2n-1)}6+(n^2-n+1)n^2
=n^4-\frac23n^3+\frac12n^2+\frac16n,
\]
\[
P_{n^2}=1-\frac{2}{3n}+\frac{1}{2n^2}+\frac{1}{6n^3}.
\]
再令 $q=\lceil\sqrt n\rceil-1$，于是 $q<\sqrt n$，且无实根只可能发生在 $1\leq a\leq q$。无实根数对数为 $\sum_{a=1}^{q}(n-a^2)$，它严格小于 $qn<n\sqrt n$。除以 $n^2$，即得 $1-P_n<1/\sqrt n$。
\end{proof}
\textbf{复盘。}这是有序、独立、允许重复的两次选取，分母是 $n^2$，不是 $\binom n2$。固定一个坐标逐列计数通常比画出全部整数点更稳妥；证明概率下界时则改数补集。

\subsection{2010 年：把有理性封闭在递推关系中}\label{sec:g2010}
\textbf{来源。}见来源~\ref{src:g2010}。
\begin{example}[题意重述]
三角形 $ABC$ 的三边长均为有理数。证明 $\cos A$ 为有理数，并证明对每个正整数 $n$，$\cos(nA)$ 也为有理数。
\end{example}
\textbf{思路。}第一问由余弦定理连接边与角；第二问需要避免引入无法保证有理的 $\sin A$。把两个余弦和角公式相加，便能消去正弦项。
\begin{proof}
设三边 $a,b,c$ 分别对着 $A,B,C$。由 $b,c>0$ 和余弦定理，
\[
t:=\cos A=\frac{b^2+c^2-a^2}{2bc}\in\mathbb Q.
\]
令 $u_n=\cos(nA)$，并补充 $u_0=1$。由和差角公式得
\begin{equation}\label{eq:cos-rec}
u_{n+1}=2tu_n-u_{n-1}\qquad(n\geq1).
\end{equation}
初值 $u_0,u_1$ 均有理。如果相邻两项 $u_{n-1},u_n$ 有理，则右端由有理数的加减乘运算构成，$u_{n+1}$ 仍有理。数学归纳法遂给出所有 $u_n\in\mathbb Q$。
\end{proof}
\textbf{复盘。}三边有理并不推出 $\sin A$ 有理，例如边长均为 $1$ 的等边三角形。递推式~\eqref{eq:cos-rec} 还表明 $u_n$ 是 $t$ 的整系数多项式；这种保持在同一数域中的思路比逐次展开倍角公式更通用。

\subsection{2011 年：差为 3 的倍数与同余分类}\label{sec:g2011}
\textbf{来源。}见来源~\ref{src:g2011}。
\begin{example}[题意重述]
设整数 $n\geq4$，点 $P(a,b)$ 的坐标满足 $a,b\in\{1,\ldots,n\}$ 且 $a>b$。求满足 $a-b=3$ 的点数 $A_n$，以及满足 $(a-b)/3$ 为整数的点数 $B_n$。
\end{example}
\textbf{思路。}第一问固定差；第二问既可枚举差 $3,6,9,\ldots$，也可把整数按除以 $3$ 的余数分组。后一种方法自动处理 $a>b$：同组任取两个不同数，大的放在第一个坐标，便唯一确定一个点。
\begin{proof}
$a=b+3$ 时，$b$ 依次取 $1,\ldots,n-3$，故 $A_n=n-3$。一般地，令 $q=\lfloor(n-1)/3\rfloor$，固定 $a-b=3j$ 得 $n-3j$ 个点，所以
\[
B_n=\sum_{j=1}^{q}(n-3j)=qn-\frac{3q(q+1)}2.
\]
若改按余数分类，令 $n=3m+r$，$r\in\{0,1,2\}$，各类大小分别为 $m,m,m$；$m+1,m,m$；$m+1,m+1,m$。于是
\[
B_n=\begin{cases}
3m(m-1)/2,&n=3m,\\
m(3m-1)/2,&n=3m+1,\\
m(3m+1)/2,&n=3m+2.
\end{cases}
\]
两种表达式计数同一集合，因而相等。
\end{proof}
\textbf{复盘。}“同余”不意味着允许两个坐标相等。用 $\binom{s}{2}$ 选同组两数，已经同时排除了对角线并消除了坐标反向的重复，不应再除以 $2$。

\subsection{2012 年：倍增链上的交替选择}\label{sec:g2012}
\textbf{来源。}见来源~\ref{src:g2012}。
\begin{example}[题意重述]
设 $P_n=\{1,2,\ldots,n\}$。满足如下条件的子集 $A\subseteq P_n$ 的个数记为 $f(n)$：若 $x\in A$，则 $2x\notin A$；若 $x\in P_n\setminus A$，则 $2x\notin P_n\setminus A$。求 $f(4)$ 和一般的 $f(n)$。
\end{example}
\textbf{思路。}当 $2x\leq n$ 时，两个条件合在一起要求 $x$ 与 $2x$ 恰好一个属于 $A$。每个正整数都能唯一写为奇数乘以 $2$ 的幂，因此把 $P_n$ 分解成以奇数开头的倍增链。
\begin{proof}
对每个奇数 $u\leq n$，取链
\[
u,2u,4u,\ldots,2^{\ell}u,
\qquad 2^{\ell}u\leq n<2^{\ell+1}u.
\]
不断除去因子 $2$，可知每个整数属于且仅属于其中一条链。在一条链内，选定首项是否属于 $A$ 后，其余项的归属必须交替，恰有两种选择。不同链之间不存在倍增约束，选择相互独立。链的条数等于不超过 $n$ 的奇数个数 $\lceil n/2\rceil$，故
\[
\boxed{f(n)=2^{\lceil n/2\rceil}}.
\]
$n=4$ 时两条链为 $1,2,4$ 和 $3$，四个集合分别为 $\{1,4\}$、$\{1,3,4\}$、$\{2\}$、$\{2,3\}$，所以 $f(4)=4$。
\end{proof}
\textbf{复盘。}若 $2x>n$，它本来就在两个集合之外，不再施加新的限制。只使用第一个条件会误算成“不同时选取相邻链点”的问题，得到完全不同的计数。这里本质是若干条路径的二染色，每条路径只有两种。

\subsection{2013 年：分组数列中的整除性}\label{sec:g2013}
\textbf{来源。}见来源~\ref{src:g2013}。
\begin{example}[题意重述]
把数列依次分组：第 $k$ 组有 $k$ 项，每项都为 $(-1)^{k-1}k$，即 $1;-2,-2;3,3,3;\ldots$。设 $S_n$ 为前 $n$ 项和，
\[
P_l=\{n\in\mathbb N^*:1\leq n\leq l,\ S_n\text{ 是 }a_n\text{ 的整数倍}\}.
\]
求 $|P_{11}|$ 与 $|P_{2000}|$。
\end{example}
\textbf{思路。}先把位置拆成“第几组、组内第几项”，而不要从 $1$ 到 $2000$ 逐项算。整除问题转化成商是否为整数后，组内位置会消失，只剩组号的奇偶性。
\begin{proof}
令 $t_k=k(k+1)/2$ 为第 $k$ 组末项的位置，$t_0=0$。交错平方和满足
\[
S_{t_k}=1^2-2^2+\cdots+(-1)^{k-1}k^2
=(-1)^{k-1}\frac{k(k+1)}2.
\]
这个等式可从 $k=1$ 出发验证：若第 $k$ 步成立，则再加上下一组的和 $(-1)^k(k+1)^2$，得到 $(-1)^k(k+1)(k+2)/2$，恰是第 $k+1$ 步的右端。若 $n=t_{k-1}+j$，其中 $1\leq j\leq k$，则
\[
\frac{S_n}{a_n}
=\frac{(-1)^{k-2}(k-1)k/2+j(-1)^{k-1}k}{(-1)^{k-1}k}
=j-\frac{k-1}{2}.
\]
因此恰好奇数组的全部位置满足条件，偶数组的全部位置不满足。$t_4=10$，前 $11$ 项取第 $1,3$ 组及第 $5$ 组第一项，共 $1+3+1=5$ 个。又 $t_{62}=1953<2000<t_{63}=2016$，故
\[
|P_{2000}|=(1+3+\cdots+61)+(2000-1953)=31^2+47=1008.
\]
\end{proof}
\textbf{复盘。}必须区分组号 $k$ 与项号 $n$；最后一组通常不完整，不能一并算入全部 $63$ 项。

\subsection{2014 年：保持形式的高阶导数关系}\label{sec:g2014}
\textbf{来源与校勘。}见来源~\ref{src:g2014}。现有纯文本转录未保留第二问左端的绝对值，下面按绝对值形式给出经推导核验的题意；这属于数学上的校勘，原卷该处符号仍需进一步核对，不能仅凭等式可成立就认定已核得原文。
\begin{example}[题意重述]
设 $f_0(x)=\sin x/x$，$x>0$，且 $f_n=f_{n-1}'$，$n\in\mathbb N^*$。求 $2f_1(\pi/2)+(\pi/2)f_2(\pi/2)$，并证明
\[
\left|n f_{n-1}(\pi/4)+\frac\pi4f_n(\pi/4)\right|=\frac{\sqrt2}{2}\qquad(n\geq1).
\]
\end{example}
\textbf{思路。}所求并非单独的 $f_n$，而是两个相邻导数的特殊线性组合。先把分母移走，得到 $xf_0(x)=\sin x$，求导后恰好产生这样的组合。
\begin{proof}
对一切正整数 $n$，都有
\begin{equation}\label{eq:g2014}
xf_n(x)+nf_{n-1}(x)=\sin\left(x+\frac{n\pi}{2}\right).
\end{equation}
$n=1$ 时由 $xf_0=\sin x$ 求导成立。假设第 $n$ 阶等式成立，对 $x$ 求导，左端为 $xf_{n+1}+(n+1)f_n$，右端为 $\sin(x+(n+1)\pi/2)$，归纳完成。所有函数在 $x>0$ 上任意阶可导，故每一步都有依据。

第一问令 $n=2,x=\pi/2$，得结果 $\sin(3\pi/2)=-1$。第二问令 $x=\pi/4$，由 $\sin(\pi/4+n\pi/2)$ 的值只能为 $\pm\sqrt2/2$，取绝对值即得结论。
\end{proof}
\textbf{复盘。}若直接展开 $f_n$，表达式会迅速膨胀。观察所求式能否由一个简单乘积连续求导获得，往往比求出每个导数的通项更有效。若删去绝对值，$n=2$ 已给出反例。

\subsection{2015 年：整除关系、取整函数与数学归纳法}\label{sec:g2015}
\textbf{来源。}见来源~\ref{src:g2015}。
\begin{example}[题意重述]
令 $X=\{1,2,3\}$、$Y_n=\{1,2,\ldots,n\}$，并设
\[
S_n=\{(a,b):a\in X,\ b\in Y_n,\ a\mid b\text{ 或 }b\mid a\},\qquad f(n)=|S_n|.
\]
求 $f(6)$，给出 $n\geq6$ 时的 $f(n)$，并用数学归纳法证明。
\end{example}
\textbf{思路。}先按 $a=1,2,3$ 分类。这三个集合互不相交，而每一类内部的“或”不能重复计数。例如 $a=b=2$ 只算一次。取整函数可以避免冗长的模 $6$ 分段。
\begin{proof}
$a=1$ 时任意 $b$ 都可以，共 $n$ 个；$a=2$ 时允许所有偶数 $b$ 以及 $b=1$，共 $\lfloor n/2\rfloor+1$ 个；$a=3$ 时允许所有 $3$ 的倍数以及 $b=1$，共 $\lfloor n/3\rfloor+1$ 个。因此猜想并得到
\begin{equation}\label{eq:g2015}
f(n)=n+\left\lfloor\frac n2\right\rfloor+\left\lfloor\frac n3\right\rfloor+2,
\qquad f(6)=6+3+2+2=13.
\end{equation}
下面另按题目要求完成归纳证明。$n=6$ 已成立。设 $n=k\geq6$ 时成立，扩展到 $k+1$，只会增加第二坐标为 $k+1$ 的数对。由于 $k+1>3$，不可能有 $k+1\mid a$；而 $a\mid k+1$ 对 $a=1$ 恒成立，对 $a=2,3$ 分别由相应整除条件决定。故
\[
f(k+1)-f(k)=1+\boldsymbol1_{2\mid k+1}+\boldsymbol1_{3\mid k+1}.
\]
其中 $\boldsymbol1_E$ 为命题 $E$ 成立时取 $1$、否则取 $0$ 的示性数。因为
\[
\left\lfloor\frac{k+1}{d}\right\rfloor-\left\lfloor\frac{k}{d}\right\rfloor
=\boldsymbol1_{d\mid k+1}\qquad(d=2,3),
\]
式~\eqref{eq:g2015} 的右端也增加同样的量，所以归纳步成立。
\end{proof}
\textbf{复盘。}分类计数已能直接证明通式，但题目明确要求数学归纳法时还应写出初始步与新增数对，不能仅写“易由归纳法得证”。

\subsection{2016 年：消去权因子，化为组合数逐项累加}\label{sec:g2016}
\textbf{来源。}见来源~\ref{src:g2016}。
\begin{example}[题意重述]
计算 $7\binom63-4\binom74$；若 $m,n\in\mathbb N^*$ 且 $n\geq m$，证明
\[
\sum_{k=m}^{n}(k+1)\binom km=(m+1)\binom{n+2}{m+2}.
\]
\end{example}
\textbf{思路。}第一问不只是数值计算，也暗示 $(k+1)\binom km=(m+1)\binom{k+1}{m+1}$。权 $k+1$ 被转移成与求和变量无关的 $m+1$ 后，可以使用式~\eqref{eq:hockey}。
\begin{proof}
第一问结果为 $7\cdot20-4\cdot35=0$。由阶乘定义及引理~\ref{lem:binomial}，
\[
\sum_{k=m}^{n}(k+1)\binom km
=(m+1)\sum_{k=m}^{n}\binom{k+1}{m+1}
=(m+1)\binom{n+2}{m+2}.
\]
若不用已知求和公式，则将 $\binom{k+1}{m+1}$ 写成 $\binom{k+2}{m+2}-\binom{k+1}{m+2}$，相加后首端项 $\binom{m+1}{m+2}=0$，只留下末端项。
\end{proof}
\textbf{复盘。}求和下限对应的组合数为零，是裂项后没有额外常数项的原因；可代入边界 $n=m$ 检查右端下标有没有错移。

\subsection{2017 年：最后一次出现的位置与倒数期望}\label{sec:g2017}
\textbf{来源。}见来源~\ref{src:g2017}。
\begin{example}[题意重述]
袋中有 $m$ 个白球、$n$ 个黑球，$m,n\in\mathbb N^*$、$n\geq2$。无放回地随机依次取球，依次放入编号 $1,2,\ldots,m+n$ 的抽屉。求第 $2$ 个抽屉放入黑球的概率。若 $X$ 是最后一个黑球所在抽屉编号的倒数，证明
\[
\mathbb E X<\frac{n}{(m+n)(n-1)}.
\]
\end{example}
\textbf{思路。}把随机过程改写成“在全部抽屉位置中等可能地选择 $n$ 个黑球位置”。最后位置为 $k$ 时，前 $k-1$ 个位置恰放另 $n-1$ 个黑球；期望求和中的 $1/k$ 与组合数不能直接相消，但将它放大为 $1/(k-1)$ 后便能化简。
\begin{proof}
记 $N=m+n$。每个位置成为黑球位置的概率相同，故第一问为 $n/N$。也可直接计算 $\binom{N-1}{n-1}/\binom Nn=n/N$。

令最后一个黑球的位置为 $K$，则 $K\in\{n,n+1,\ldots,N\}$，$X=1/K$，并且
\[
\mathbb P(K=k)=\frac{\binom{k-1}{n-1}}{\binom Nn}.
\]
由于 $k\geq n\geq2$，$1/k<1/(k-1)$，且各项概率均为正。因此
\begin{align*}
\mathbb E X
&=\frac1{\binom Nn}\sum_{k=n}^{N}\frac1k\binom{k-1}{n-1}\\
&<\frac1{(n-1)\binom Nn}\sum_{k=n}^{N}\binom{k-2}{n-2}
=\frac{\binom{N-1}{n-1}}{(n-1)\binom Nn}
=\frac n{N(n-1)}.
\end{align*}
中间等式使用了式~\eqref{eq:hockey}。
\end{proof}
\textbf{复盘。}无放回抽取的颜色序列相互依赖，不是独立的伯努利试验；但黑球位置集合仍等可能。另有 $\mathbb E(1/K)\ne1/\mathbb E K$，不能用求平均位置替代本题。严格不等号在逐项放大时已经得到，不需要额外寻找不能取等的末尾情形。

\subsection{2018 年：逆序数为 2 的排列}\label{sec:g2018}
\textbf{来源。}见来源~\ref{src:g2018}。
\begin{example}[题意重述]
把 $1,2,\ldots,n$ 排列，记其中逆序数恰为 $k$ 的排列个数为 $f_n(k)$。求 $f_3(2)$、$f_4(2)$，并求 $n\geq5$ 时的 $f_n(2)$。
\end{example}
\textbf{思路。}插入最大元素的位置直接决定新增逆序数。目标只有 $2$，所以最大元素必须位于末尾、倒数第二、倒数第三三个位置之一。这把复杂的排列枚举化成三个规模更小的问题。
\begin{proof}
由引理~\ref{lem:insert}，$f_n(0)=1$，且
\[
f_n(1)=f_{n-1}(1)+f_{n-1}(0)=f_{n-1}(1)+1.
\]
结合 $f_2(1)=1$，得 $f_n(1)=n-1$。对 $n\geq3$，
\[
f_n(2)=f_{n-1}(2)+f_{n-1}(1)+f_{n-1}(0)
=f_{n-1}(2)+n-1.
\]
由于 $f_2(2)=0$，故
\[
f_n(2)=\sum_{j=3}^{n}(j-1)=\frac{n(n-1)}2-1
=\frac{(n-2)(n+1)}2.
\]
特别地 $f_3(2)=2$，对应排列 $231,312$；$f_4(2)=5$。
\end{proof}
\textbf{复盘。}递推必须覆盖新增逆序数为 $0,1,2$ 三种互斥情形；最大元素放得更靠前会造成至少 $3$ 个逆序，不能计入。公式还适用于 $n=2,3,4$，但不能把 $n=1$ 代进去当作有效计数。

\subsection{2019 年：三条水平线上的格点距离}\label{sec:g2019}
\textbf{来源。}见来源~\ref{src:g2019}。
\begin{example}[题意重述]
设
\[
A_n=\{(i,0):i=0,1,\ldots,n\},\quad
B_n=\{(0,1),(n,1)\},\quad
C_n=\{(i,2):i=0,1,\ldots,n\},
\]
$D_n=A_n\cup B_n\cup C_n$。从 $D_n$ 中等可能地选两个不同点，距离记为随机变量 $X$。求 $n=1$ 时 $X$ 的分布列；对 $n\geq3$，求 $\mathbb P(X\leq n)$。
\end{example}
\textbf{思路。}小规模直接按平方距离分类；一般规模改数距离超过 $n$ 的少数点对。横坐标差至多 $n$，若不到 $n$，纵坐标差至多 $2$，距离平方就能被统一控制。
\begin{proof}
$n=1$ 时共有六个点、$\binom62=15$ 个无序点对。平方距离 $1,2,4,5$ 对应的点对数依次为 $7,4,2,2$，所以
\[
\begin{array}{c|cccc}
X&1&\sqrt2&2&\sqrt5\\\hline
\mathbb P&7/15&4/15&2/15&2/15
\end{array}
\]
例如距离 $1$ 包括三条水平单位线段与四条竖直单位线段；距离 $\sqrt2$ 为两个单位方格的四条对角线，其余四对跨越上下两行。

一般情形共有 $2n+4$ 个点。设一对点的横、纵坐标差绝对值分别为 $d,e$。若 $d\leq n-1$，则
\[
d^2+e^2\leq(n-1)^2+4=n^2-2n+5<n^2\quad(n\geq3).
\]
故 $X>n$ 必须 $d=n$，即两个点分别位于 $x=0$ 与 $x=n$。两条竖线上各有三个点，共 $9$ 对，其中同高的 $3$ 对距离等于 $n$，不同高的 $6$ 对距离大于 $n$。因此
\[
\boxed{\mathbb P(X\leq n)=1-\frac6{\binom{2n+4}{2}}
=1-\frac6{(n+2)(2n+3)}}.
\]
\end{proof}
\textbf{复盘。}“距离不超过”包含等号，应只删去不同高的极端点对。$n\geq3$ 正是统一估计成立的条件；$n=2$ 时 $(d,e)=(1,2)$ 也给出超过阈值的距离，不能沿用同一补集计数。

\section{市级模拟与调研卷：已核题解}\label{sec:city}
本节按年份排列；同一联考仅列一次。来源信息保留考试名称和转载页，题干用本书的统一符号重新叙述，解答重新组织推导。页码均指来源网站的文档分页，不是本书页码。完整检索情况见第~\ref{sec:register} 节。

\subsection{2010 年南京一模：准线上的乘积条件与定点弦}\label{sec:c2010nj1}
\textbf{来源。}\href{https://doc.xuehai.net/b9a0d69dda6b154ccecee1997-11.html}{南京市 2010 届高三第一次模拟考试，第 11 页}，附加题第 6 题；答案在同文档后续页。另有标题写作“二模”的重复转载，暂不据此增加一份试卷。
\begin{example}[题意]
抛物线顶点为 $O(0,0)$，焦点为 $(1,0)$。准线上两点 $M,N$ 的纵坐标分别为 $u,v$，且 $uv=-4$。直线 $OM,ON$ 分别与抛物线再交于 $A,B$。求抛物线方程，并证明直线 $AB$ 经过定点。
\end{example}
\textbf{思路。}准线条件控制的是两个参数的乘积。把抛物线上的点写成 $(s^2,2s)$，弦方程中的常数项正好含有两参数之积，可以直接找到定点。
\begin{proof}
由焦点坐标得抛物线为 $y^2=4x$，准线为 $x=-1$。因 $uv=-4$，两数均非零。直线 $OM$ 为 $y=-ux$，与抛物线联立，除 $O$ 外的交点为
\[
A=\left(\frac4{u^2},-\frac4u\right)=(s^2,2s),\qquad s=-\frac2u.
\]
同理 $B=(t^2,2t)$，其中 $t=-2/v$，故 $st=4/(uv)=-1$。特别地 $s,t$ 异号，$A\ne B$。将两点代入可验证其连线方程为
\[
2x-(s+t)y+2st=0,
\]
即 $2x-(s+t)y-2=0$，所以直线恒过 $(1,0)$。
\end{proof}
\textbf{复盘。}用一般式写弦方程可以包含 $s+t=0$ 时的竖直弦；若先除以 $s+t$ 求斜率，就必须另补这一情形。本题定点恰为焦点，而这一结论来自 $st=-1$，不需要凭图形猜测。

\subsection{2010 年南通、扬州、泰州二模：递推轨道的有界性}\label{sec:c2010nty2}
\textbf{来源。}\href{https://doc.xuehai.net/b0564b8cf950bcb418f9af425-5.html}{三市 2010 届高三第二次模拟考试，第 5--6 页}，第 23 题；同文档第 14--15 页为解答。
\begin{example}[题意]
给定实数 $a$，定义 $a_1=a$、$a_{n+1}=a_n^2+a$，并令
\[
M=\{a\in\mathbb R:\text{对每个正整数 }n,\ |a_n|\leq2\}.
\]
证明：$a<-2$ 时 $a\notin M$；$0<a\leq1/4$ 时 $a\in M$；判断 $a>1/4$ 时是否属于 $M$ 并证明。
\end{example}
\textbf{思路。}第一种情形看首项即可。第二种情形寻找被递推映射保持的区间。第三种情形不能仅说“递增所以无界”，应找出每一步增长量的统一正下界。
\begin{proof}
若 $a<-2$，则 $|a_1|>2$，直接违背定义。若 $0<a\leq1/4$，归纳证明 $0<a_n\leq1/2$：首项满足，且当 $0<a_n\leq1/2$ 时，
\[
0<a_{n+1}=a_n^2+a\leq\frac14+\frac14=\frac12.
\]
因此所有项的绝对值都不超过 $2$，即 $a\in M$。

若 $a>1/4$，设 $\varepsilon=a-1/4>0$。则
\[
a_{n+1}-a_n=\left(a_n-\frac12\right)^2+\varepsilon\geq\varepsilon.
\]
逐项相加得 $a_n\geq a+(n-1)\varepsilon$。取足够大的正整数 $n$ 使右端大于 $2$，便有 $a_n>2$，故 $a\notin M$。
\end{proof}
\textbf{复盘。}“对所有 $n$”是集合定义的关键量词。证明属于 $M$ 需要统一控制全部项，证明不属于只需找到一个越界项。原题这三问没有要求确定整个 $M$，不应把未讨论的 $[-2,0]$ 擅自排除。

\subsection{2011 年南京、盐城三模：正四面体游走与二项式加权}\label{sec:c2011ny3}
\textbf{来源。}\href{https://doc.xuehai.net/b699d388c622547cec887cbae-6.html}{南京、盐城 2011 届高三第三次模拟考试，第 6 页}，第 23 题。
\begin{example}[题意]
小虫从正四面体 $ABCD$ 的顶点 $A$ 出发，每一步沿一条棱等概率地到达另一个顶点。令 $p_n$ 为第 $n$ 步后位于 $A$ 的概率。求 $p_1,p_2,p_3$ 及 $p_n$ 的表达式；再求
\[
S_n=\sum_{j=1}^{n}\binom njp_j.
\]
\end{example}
\textbf{思路。}从 $A$ 必然离开，从其余任一顶点返回 $A$ 的概率都是 $1/3$。所以“是否在 $A$”已经包含递推所需的全部状态信息。解出递推后，加权和就是在两个点取值的二项式展开式。
\begin{proof}
设 $p_0=1$，则 $p_{n+1}=(1-p_n)/3$，依次得到 $p_1=0$、$p_2=1/3$、$p_3=2/9$。又有
\[
p_{n+1}-\frac14=-\frac13\left(p_n-\frac14\right),
\]
故 $p_n=1/4+(3/4)(-1/3)^n$。于是
\begin{align*}
S_n
&=\frac14\sum_{j=1}^{n}\binom nj
 +\frac34\sum_{j=1}^{n}\binom nj\left(-\frac13\right)^j\\
&=\frac14(2^n-1)+\frac34\left[\left(\frac23\right)^n-1\right]\\
&=\boxed{2^{n-2}+\frac34\left(\frac23\right)^n-1}.
\end{align*}
\end{proof}
\textbf{复盘。}求和从 $1$ 开始，两次使用二项式定理都要减去 $j=0$ 项。与第~\ref{sec:c2013ny3} 节相比，这里的递推系数为负，因此概率围绕 $1/4$ 交替波动；两个模型的状态转移不能混用。

\subsection{2012 年南通、泰州、扬州一调：两种交点共享横坐标}\label{sec:c2012nty1}
\textbf{来源。}\href{https://doc.xuehai.net/b8780209d98e4d302aafd63ce-6.html}{三市 2012 届高三第一次调研测试，第 6 页}，第 23 题，后续页含答案。
\begin{example}[题意]
抛物线顶点为原点，焦点为 $(1,0)$。其上不同两点 $A,B$ 均在 $x$ 轴上方，过它们的切线分别交 $x$ 轴于 $C,D$，两切线交于 $M$，$AD$ 与 $BC$ 交于 $N$。求抛物线方程，证明 $MN\perp x$ 轴；若 $MN$ 与 $x$ 轴相交于 $(1,0)$，证明 $AB$ 过定点。
\end{example}
\textbf{思路。}把各点坐标都化为两个切点参数，就能发现 $M,N$ 的横坐标相同。定点问题随后转化为参数乘积固定。
\begin{proof}
抛物线为 $y^2=4x$。写 $A=(s^2,2s)$、$B=(t^2,2t)$，其中 $s,t>0$ 且 $s\ne t$。切线分别为 $sy=x+s^2$、$ty=x+t^2$，所以
\[
C=(-s^2,0),\quad D=(-t^2,0),\quad M=(st,s+t).
\]
直线 $AD,BC$ 的方程分别为
\[
y=\frac{2s(x+t^2)}{s^2+t^2},\qquad
y=\frac{2t(x+s^2)}{s^2+t^2}.
\]
联立并利用 $s\ne t$，得
\[
N=\left(st,\frac{2st(s+t)}{s^2+t^2}\right).
\]
由于 $s^2+t^2>2st$ 且 $s+t>0$，$N$ 的纵坐标严格小于 $M$ 的纵坐标。因此 $M\ne N$，两点确定的直线就是 $x=st$，从而垂直于 $x$ 轴。

若这条直线通过 $(1,0)$，则 $st=1$。弦 $AB$ 的方程 $2x-(s+t)y+2st=0$ 化为 $2x-(s+t)y+2=0$，恒过 $(-1,0)$。
\end{proof}
\textbf{复盘。}证明“横坐标相同”之后还应验证 $M\ne N$，否则无法谈论直线 $MN$。条件“两点在 $x$ 轴上方”在这里保证参数为正，也排除了交点退化。

\subsection{2013 年苏锡常镇一调：二项分布的绝对偏差}\label{sec:c2013sxcz1}
\textbf{来源。}\href{https://doc.xuehai.net/b0d470082d90b0ed78020c4e5-6.html}{苏锡常镇四市 2013 届高三教学情况调研（一），第 6 页}，第 23 题。
\begin{example}[题意]
某人对每一景点是否游览相互独立，游览概率都是 $1/2$。设 $\xi$ 为游览景点数与未游览景点数之差的绝对值。先在共有 $3$ 个景点时求 $\xi$ 的分布与期望，再对奇数 $n\geq3$ 个景点求一般结果。
\end{example}
\textbf{思路。}先记录游览次数 $Y$，再把 $\xi=|2Y-n|$ 看作 $Y$ 的函数。两侧对称的 $Y$ 值映到同一个 $\xi$ 值，概率必须相加；期望则可借相邻组合数之差消去求和。
\begin{proof}
有 $Y\sim\mathrm{Bin}(n,1/2)$。当 $n=3$ 时，$\xi$ 只取 $1,3$，且
\[
\mathbb P(\xi=1)=\frac34,\qquad
\mathbb P(\xi=3)=\frac14,\qquad \mathbb E\xi=\frac32.
\]
一般地，令 $n=2r+1$。对 $0\leq j\leq r$，$|2Y-n|=2j+1$ 对应 $Y=r-j$ 或 $Y=r+j+1$。由组合数对称性，
\[
\mathbb P(\xi=2j+1)=\frac{2\binom n{r-j}}{2^n}.
\]
利用关于 $n/2$ 的对称性，只需计算正半边：
\[
\mathbb E\xi=\frac2{2^n}\sum_{k=r+1}^{n}(2k-n)\binom nk.
\]
又因为
\[
(2k-n)\binom nk
=n\left[\binom{n-1}{k-1}-\binom{n-1}k\right],
\]
相邻项相消，故
\[
\boxed{\mathbb E\xi=\frac{n\binom{n-1}r}{2^{n-1}}.}
\]
\end{proof}
\textbf{复盘。}$\mathbb E|2Y-n|$ 不能替换为 $|2\mathbb EY-n|$；后者为零，前者严格为正。奇数条件使 $\xi$ 不取零，因而对称配对时没有中间项需要单独处理。

\subsection{2013 年南京、盐城三模：图上游走与倒数和}\label{sec:c2013ny3}
\textbf{来源。}\href{https://doc.xuehai.net/b764d9602946662646540b4b9-6.html}{南京、盐城 2013 届高三第三次模拟考试，第 6 页}，第 23 题；第 15--16 页为解答。
\begin{example}[题意]
一个动点从三棱柱上底面的某一顶点出发，每步等概率沿相接的三条棱之一走到另一端顶点。令 $p_n$ 为走 $n$ 步后位于上底面顶点的概率。求 $p_1,p_2$，并证明对正整数 $n$ 有
\[
\sum_{i=1}^{n}\frac1{4p_i-1}>\frac{n^2}{n+1}.
\]
\end{example}
\textbf{思路。}六个顶点只需合并为上、下两个状态。先求概率递推式，再处理确定的数列不等式，可以避免同时展开大量路径。
\begin{proof}
从上底面出发，下一步仍在上底面的概率为 $2/3$；从下底面出发，下一步进入上底面的概率为 $1/3$。令 $p_0=1$，全概率公式给出
\[
p_{n+1}=\frac23p_n+\frac13(1-p_n)=\frac{p_n+1}{3}.
\]
故 $p_1=2/3$、$p_2=5/9$。递推式减去不动点 $1/2$ 后变成等比关系，得到
\[
p_n=\frac12+\frac1{2\cdot3^n},\qquad d_i:=4p_i-1=1+\frac2{3^i}>0.
\]
对正数 $d_i$，柯西不等式给出 $(\sum d_i)(\sum1/d_i)\geq n^2$。又有
\[
\sum_{i=1}^{n}d_i=n+1-3^{-n}<n+1,
\]
所以
\[
\sum_{i=1}^{n}\frac1{4p_i-1}
\geq\frac{n^2}{n+1-3^{-n}}>\frac{n^2}{n+1}.
\]
\end{proof}
\textbf{复盘。}最终的严格不等号来自有限等比和中的 $3^{-n}>0$，并不需要证明柯西不等式严格。即使 $n=1$ 时柯西等号成立，结论仍然严格。

\subsection{2014 届苏北四市期末调研：切点弦与平行条件}\label{sec:c2014sb}
\textbf{来源与名称。}\href{https://doc.xuehai.net/b24b27f9a1f1345ca980b6ebe-5.html}{徐州、淮安、连云港、宿迁四市试卷，第 5 页}，第 23 题；该转载题名用“一模”，另有\href{https://www.mianfeiwendang.com/doc/e71df5c172c3e51967676141}{期末调研转载}。本书保留“期末调研”，不把两个名称重复计数。
\begin{example}[题意]
给定 $A=(-1,0)$、$F=(1,0)$，动点 $P$ 满足 $\overrightarrow{AP}\cdot\overrightarrow{AF}=2|FP|$。求其轨迹 $C$。在直线 $l:y=2x+2$ 上取点 $Q$，由 $Q$ 向 $C$ 作两条不同切线，切点为 $M,N$。是否存在这样的 $Q$，使 $MN\parallel l$？
\end{example}
\textbf{思路。}向量条件是点到焦点距离与到准线距离相等的代数形式。第二问先写一般切点弦，再施加平行条件；最后核查确实有两个不同实切点。
\begin{proof}
设 $P=(x,y)$。已知条件为 $x+1=\sqrt{(x-1)^2+y^2}$，平方得到 $y^2=4x$。反过来，若 $y^2=4x$，则 $x\geq0$，从而 $x+1>0$，可以还原原等式。因此轨迹恰为 $C:y^2=4x$。

设 $Q=(u,v)$，切点用 $(t^2,2t)$ 表示，切线为 $ty=x+t^2$。它经过 $Q$ 当且仅当
\[
t^2-vt+u=0.
\]
设两个不同实根为 $s,t$，则 $s+t=v$、$st=u$，切点弦为 $2x-vy+2u=0$。若与斜率为 $2$ 的 $l$ 平行，则 $v\ne0$ 且 $2/v=2$，即 $v=1$。由 $Q\in l$ 得 $u=-1/2$。

对 $Q=(-1/2,1)$，参数方程的判别式为 $v^2-4u=3>0$，所以确有两个不同切点；切点弦为 $y=2x-1$，与 $l$ 平行。因此所求点存在且唯一，为 $(-1/2,1)$。
\end{proof}
\textbf{复盘。}网页转录可能漏掉最终横坐标的负号；代回 $y=2x+2$ 即可核实。证明存在性不能停在解出候选坐标，还应检查切线的实数性、不同性和两直线不重合。

\subsection{2014 年南京、盐城一模：区间和限制下的交替符号}\label{sec:c2014ny1}
\textbf{来源。}\href{https://www.51wendang.com/doc/b354a966ae74314766a3ac38}{南京、盐城 2014 届高三第一次模拟考试}，第 4--5 页第 23 题，第 11 页解答。
\begin{example}[题意]
给定正整数 $m$，考虑每项均为 $2$ 或 $-2$ 的有序数组 $(a_1,\ldots,a_{2m})$。求每一对 $(a_{2k-1},a_{2k})$ 都异号的数组数 $A$。再要求任意 $1\leq k\leq l\leq m$ 都满足
\[
\left|\sum_{i=2k-1}^{2l}a_i\right|\leq4,
\]
求其中至少有一对 $(a_{2k-1},a_{2k})$ 同号的数组数 $B$。
\end{example}
\textbf{思路。}把两个数看成一组，缩放后的组和只可能为 $-1,0,1$。所有连续区间和的限制，等价于删去零后正负号交替；零组还保留两种内部排列，不能少算。
\begin{proof}
异号组有 $(2,-2),(-2,2)$ 两种，各组独立选择，因此 $A=2^m$。令 $b_k=(a_{2k-1}+a_{2k})/4$，则 $b_k\in\{-1,0,1\}$，且题设等价于每一段连续 $b_k$ 的和都在 $[-1,1]$ 内。

删去所有零项后，相邻的两个非零项必须异号：否则它们以及中间的零组成的连续区间和为 $2$ 或 $-2$。反过来，若非零项交替，则任意连续区间含有的非零项也是一段交替序列，其和只可能为 $-1,0,1$。所以这一条件既必要又充分。

若恰有 $j\geq1$ 个非零组，选择位置有 $\binom mj$ 种；第一非零项有两种符号，之后符号全部确定；每个零组有两种内部排列。故
\[
B=\sum_{j=1}^{m}2\binom mj2^{m-j}
=2\left[(2+1)^m-2^m\right]
=\boxed{2(3^m-2^m)}.
\]
\end{proof}
\textbf{复盘。}这里限制的是所有连续的“整组区间”，不是所有单项区间。若错误地让起止下标任意取值，条件就变强了。计数前先证明结构条件的充分性，可防止只数到一个必要条件所允许的更大集合。

\subsection{2015 年南京、盐城一模：有序分离的两个子集}\label{sec:c2015ny1}
\textbf{来源。}\href{https://doc.xuehai.net/b440b636eecd6729ce54c837f-8.html}{南京、盐城 2015 届高三第一次模拟考试，第 8 页}，第 23 题及解答。
\begin{example}[题意]
设整数 $n\geq2$，从 $\{1,2,\ldots,n\}$ 中选两个非空子集 $A,B$，满足 $\max A<\min B$，这样的有序对 $(A,B)$ 的数目记为 $P_n$。求 $P_2,P_3$，并求一般的 $P_n$。
\end{example}
\textbf{思路。}直接枚举每个元素属于 $A$、$B$ 或两者都不属于，会让顺序限制纠缠在一起。固定左侧子集的最大元后，其余选择变得独立。
\begin{proof}
当 $n=2$，仅有 $(\{1\},\{2\})$，故 $P_2=1$。一般地，固定 $\max A=k$，其中 $1\leq k\leq n-1$。元素 $k$ 必须选入 $A$，小于 $k$ 的元素任意选，故 $A$ 有 $2^{k-1}$ 种；$B$ 是 $\{k+1,\ldots,n\}$ 的非空子集，有 $2^{n-k}-1$ 种。各类按唯一的最大元区分，互不重叠，因此
\begin{align*}
P_n&=\sum_{k=1}^{n-1}2^{k-1}(2^{n-k}-1)\\
&=(n-1)2^{n-1}-(2^{n-1}-1)
=\boxed{(n-2)2^{n-1}+1}.
\end{align*}
代入 $n=3$ 得 $P_3=5$。
\end{proof}
\textbf{另一条思路。}若 $|A\cup B|=j\geq2$，先选并集，再把其递增排列在相邻两元素之间切开，共 $j-1$ 个切口，左段为 $A$、右段为 $B$。故 $P_n=\sum_{j=2}^{n}(j-1)\binom nj$，由二项式及加权二项式和也可得到同一结果。两种计数分别固定“边界元素”和“整体规模”，可相互校验。

\subsection{2015 年泰州二模：多项式恒等式与系数对称}\label{sec:c2015tz2}
\textbf{来源。}\href{https://www.51wendang.com/doc/ade6c3141d94fc4b359d414c/6}{泰州市 2015 届高三第二次模拟考试，第 6 页}，第 23 题；第 13--15 页含解答。原函数中的两个连接符均为加号。
\begin{example}[题意]
设正整数 $n$，$f(x)=(x^2+x+1)^n$。若 $2n$ 次多项式 $g$ 满足对每个实数 $x$ 都有 $f(x^2)g(x)=g(x^3)$，求 $g(1),g(-1)$，并给出一个这样的 $g$。再证明存在常数 $a_0,\ldots,a_n$ 使
\[
f(x)=\sum_{j=0}^{n-1}a_j(x^j+x^{2n-j})+a_nx^n.
\]
\end{example}
\textbf{思路。}第一问先在 $x=\pm1$ 代入，再利用 $x^6-1=(x^2-1)(x^4+x^2+1)$ 构造解。第二问需要的不是展开每个系数，而是把 $x$ 与 $1/x$ 互换后发现系数反向不变。
\begin{proof}
因为 $f(1)=3^n\ne1$，代入 $x=1,-1$ 分别得到 $(3^n-1)g(1)=0$、$(3^n-1)g(-1)=0$，所以两值均为零。取
\[
g(x)=(x^2-1)^n,
\]
则次数恰为 $2n$，且
\[
f(x^2)g(x)=[(x^4+x^2+1)(x^2-1)]^n
=(x^6-1)^n=g(x^3),
\]
符合要求。

写 $f(x)=\sum_{j=0}^{2n}c_jx^j$。对 $x\ne0$，有
\[
x^{2n}f(1/x)=(1+x+x^2)^n=f(x).
\]
两边都是多项式，在无限多个实数处相等，故是多项式恒等式。比较系数得 $c_j=c_{2n-j}$。令 $a_j=c_j$，将两端对称位置配对，中间项单列，便得到所求表达式。
\end{proof}
\textbf{复盘。}题目只要求给出一个 $g$，不要求刻画所有解。第二问中 $x=0$ 不能直接代入 $f(1/x)$，但多项式恒等原理会把结论延伸到 $0$，因而不存在定义域缺口。

\subsection{2016 年南通、扬州、泰州二模：模四余数的对称计数}\label{sec:c2016nty2}
\textbf{来源。}\href{https://www.renrendoc.com/paper/165058812.html}{三市 2016 届高三第二次模拟考试转载}，附加题末题在该题库重排为第 24 题。
\begin{example}[题意]
给定正整数 $k$，每个 $a_i\in\{0,1\}$，$1\leq i\leq4k$。令 $m(b)$ 为满足 $\sum_{i=1}^{4k}a_i\equiv b\pmod4$ 的有序数组数。求 $k=2$ 时的 $m(1)$，并证明一般情况下 $m(3)=2^{4k-2}$。
\end{example}
\textbf{思路。}先把余数 $1,3$ 合并为“奇数和”，再用逐位取反把这两类一一对应。这比直接计算相隔四项的组合数和更短。
\begin{proof}
当 $k=2$，数组有 $8$ 项，总和余 $1$ 只能取 $1,5$，所以 $m(1)=\binom81+\binom85=64$。

一般地，翻转第一位 $a_1$ 而保留其余各位，会在奇数和数组与偶数和数组之间建立双射，所以奇数和数组共有 $2^{4k-1}$ 个，即 $m(1)+m(3)=2^{4k-1}$。再将每一位都替换为 $1-a_i$，总和 $S$ 变为 $4k-S$，模 $4$ 的余数 $1$ 与 $3$ 互换。这也是可逆操作，故 $m(1)=m(3)$，得到
\[
\boxed{m(3)=2^{4k-2}}.
\]
\end{proof}
\textbf{复盘。}两次翻转承担不同任务：翻一位比较奇偶，翻全部比较模四余数。题设长度为 $4k$，正是全部取反将余数 $1$ 变成 $3$ 的原因。

\subsection{2017 年南通一模：焦点弦与三角形面积最小值}\label{sec:c2017nt1}
\textbf{来源。}\href{https://www.renrendoc.com/paper/202096368.html}{南通市 2017 届高三第一次调研测试转载}，附加题末题在该题库重排为第 26 题。
\begin{example}[题意]
抛物线 $C:x^2=2py$（$p>0$）上有点 $M=(m,1)$，它到焦点 $F$ 的距离为 $2$。求 $C$ 的方程。取 $C$ 上异于顶点的点 $E$，其切线交 $x$ 轴于 $P$，直线 $PF$ 与 $C$ 相交于 $A,B$。求 $\triangle EAB$ 面积的最小值。
\end{example}
\textbf{思路。}用焦点到准线的距离关系确定 $p$，然后只保留切点参数。面积由“弦长乘点到弦的距离”组成，可避开三个顶点坐标的行列式展开。
\begin{proof}
准线为 $y=-p/2$。由抛物线定义，$MF=1+p/2=2$，所以 $p=2$、$C:x^2=4y$、$F=(0,1)$。令 $E=(2t,t^2)$，其中 $t\ne0$。切线为 $y=tx-t^2$，所以 $P=(t,0)$，焦点弦所在直线为 $x+ty-t=0$。

代入 $y=x^2/4$，交点横坐标是 $x^2+(4/t)x-4=0$ 的两根，其差的绝对值为 $4\sqrt{1+t^2}/|t|$。弦斜率为 $-1/t$，故
\[
|AB|=\frac{4(1+t^2)}{t^2}.
\]
点 $E$ 到弦所在直线的距离为
\[
d=\frac{|2t+t^3-t|}{\sqrt{1+t^2}}
=|t|\sqrt{1+t^2}.
\]
因此 $S=2(1+t^2)^{3/2}/|t|$。令 $u=t^2>0$，由算术平均数不小于几何平均数，
\[
u+\frac12+\frac12\geq3\sqrt[3]{\frac u4},
\qquad S^2=\frac{4(1+u)^3}{u}\geq27.
\]
当且仅当 $u=1/2$ 时等号成立；此时 $E=(\pm\sqrt2,1/2)$ 确实满足条件。故最小面积为 $\boxed{3\sqrt3}$。
\end{proof}
\textbf{复盘。}$t\ne0$ 由排除顶点保证，是弦斜率和面积分母合法的前提。最小值证明既要给出下界，也要展示可达到下界的切点。

\subsection{2018 年镇江一模：分段奖励与门数差的绝对值}\label{sec:c2018zj1}
\textbf{来源与校勘。}\href{https://yn.qinxue100.com/news/gz/20180305/59236.html}{2018 镇江高三一模数学试卷及答案}，第 23 题。网页题干的“差”字表达有缺损，同页答案明确取值为 $0,2,4$，故按两类门数之差的绝对值整理。
\begin{example}[题意]
某生参加四门学科测试，各门获得 A 的概率均为 $1/4$，且相互独立。获得 A 的门数为 $Y$；当 $Y=0,1,2,3,4$ 时，奖励分 $\xi_1$ 分别为 $0,1,2,3,5$。求奖励分的数学期望；另令 $\xi_2$ 为获得 A 与未获得 A 的门数之差的绝对值，求其分布列。
\end{example}
\textbf{思路。}奖励分只在全部获得 A 时比门数多 $1$，可以直接拆成 $Y$ 与事件示性变量之和。第二个随机变量是 $|2Y-4|$，要合并对称的 $Y$ 值。
\begin{proof}
由独立性，$Y\sim\mathrm{Bin}(4,1/4)$，所以 $\mathbb EY=1$、$\mathbb P(Y=4)=1/256$。由 $\xi_1=Y+\mathbf1_{\{Y=4\}}$ 得
\[
\boxed{\mathbb E\xi_1=1+\frac1{256}=\frac{257}{256}}.
\]
$Y$ 取 $0,1,2,3,4$ 的概率依次为 $(81,108,54,12,1)/256$。故
\[
\begin{array}{c|ccc}
\xi_2&0&2&4\\\hline
\mathbb P&27/128&15/32&41/128
\end{array}
\]
例如，$\xi_2=4$ 包括 $Y=0$ 和 $Y=4$，概率为 $(81+1)/256$。三项之和为 $1$。
\end{proof}
\textbf{复盘。}四门均 A 的奖励是 $5$，不是 $4$，这一个额外单位决定期望中的 $1/256$。若把门数差错误地当成有符号差，会得到五个取值，已经改变了题目。

\subsection{2018 年盐城三模：由二元不等式到组合数倒数和}\label{sec:c2018yc3}
\textbf{来源。}\href{https://www.mianfeiwendang.com/doc/0bbc44eefbbd378979599972083302e88da33d09/6}{盐城市 2018 届高三第三次模拟考试，第 6 页}，第 23 题；第 14--15 页给出答案，可据以复核分式的位置。
\begin{example}[题意]
对正数 $a_1,a_2,b_1,b_2$，比较 $b_1^2/a_1+b_2^2/a_2$ 与 $(b_1+b_2)^2/(a_1+a_2)$，推广到任意有限项并证明。进而证明对每个正整数 $n$，
\[
\sum_{k=0}^{n}\frac{2k+1}{\binom nk}\geq\frac{(n+1)^3}{2^n}.
\]
\end{example}
\textbf{思路。}这是柯西不等式的分式形式。应用时不要直接把分子 $2k+1$ 开平方，而要写成 $(2k+1)^2/[(2k+1)\binom nk]$，使分母相加后成为可计算的组合数加权和。
\begin{proof}
直接作差，得到
\[
\frac{b_1^2}{a_1}+\frac{b_2^2}{a_2}
-\frac{(b_1+b_2)^2}{a_1+a_2}
=\frac{(a_2b_1-a_1b_2)^2}{a_1a_2(a_1+a_2)}\geq0.
\]
推广为：对 $a_i,b_i>0$，$1\leq i\leq r$，有
\begin{equation}\label{eq:engel}
\sum_{i=1}^{r}\frac{b_i^2}{a_i}\geq
\frac{(\sum_{i=1}^{r}b_i)^2}{\sum_{i=1}^{r}a_i}.
\end{equation}
其证明可逐次合并前两项，也可用数学归纳法：假设前 $r$ 项结论成立，令 $A=\sum_{i=1}^{r}a_i$、$B=\sum_{i=1}^{r}b_i$，先将前 $r$ 项合并为 $B^2/A$，再对 $B^2/A+b_{r+1}^2/a_{r+1}$ 使用已证二元式，即得 $r+1$ 项结论。单项情形为等式，归纳成立。

令 $b_{k+1}=2k+1$、$a_{k+1}=(2k+1)\binom nk$，$0\leq k\leq n$。这些数全为正，且
\[
\sum_{k=0}^{n}(2k+1)=(n+1)^2,
\qquad
\sum_{k=0}^{n}(2k+1)\binom nk
=2n2^{n-1}+2^n=(n+1)2^n.
\]
其中加权组合数和由 $k\binom nk=n\binom{n-1}{k-1}$ 得到。代入式~\eqref{eq:engel} 得
\[
\sum_{k=0}^{n}\frac{2k+1}{\binom nk}
\geq\frac{(n+1)^4}{(n+1)2^n}
=\frac{(n+1)^3}{2^n}.
\]
\end{proof}
\textbf{复盘。}不等式中分子、分母的分配决定下界能否化简。这里有意让一个奇数因子留在分母，换来了可由二项式计算的和；若简单套用 $\sum1/a_i\geq r^2/\sum a_i$，会丢失题目中的权。

\subsection{2018 年南通、泰州一模：从余弦和到加权正弦和}\label{sec:c2018nttz1}
\textbf{来源。}\href{https://yn.qinxue100.com/news/gz/20180306/59373.html}{南通、泰州市 2018 届高三一模数学试题及答案}，附加题第 23 题；该页同时保留题干和归纳、求导的解答。
\begin{example}[题意]
对正整数 $n$ 和 $x\notin2\pi\mathbb Z$，用数学归纳法证明
\[
\sum_{j=1}^{n}\cos(jx)=\frac{\sin((n+1/2)x)}{2\sin(x/2)}-\frac12.
\]
再求 $\sum_{j=1}^{2018}j\sin(j\pi/6)$。
\end{example}
\textbf{思路。}余弦求导产生负的正弦及倍角系数 $j$，所以第二问所求是第一问函数和的负导数。先把 $x$ 保留为变量，求导后再代入 $\pi/6$。
\begin{proof}
当 $n=1$ 时，由 $\sin(3x/2)-\sin(x/2)=2\sin(x/2)\cos x$，所给等式成立。假设对 $n=k$ 成立，加上 $\cos((k+1)x)$ 后，利用
\[
2\sin(x/2)\cos((k+1)x)
=\sin((k+3/2)x)-\sin((k+1/2)x),
\]
可将分子改写为 $\sin((k+3/2)x)$，得到 $n=k+1$ 时的等式，归纳完成。条件 $x\notin2\pi\mathbb Z$ 保证分母非零。

记左端为 $F_n(x)$。在分母不为零的开区间内，对有限和逐项求导，得
\[
-\sum_{j=1}^{n}j\sin(jx)
=\frac{(n+1/2)\cos((n+1/2)x)\sin(x/2)
-\tfrac12\sin((n+1/2)x)\cos(x/2)}{2\sin^2(x/2)}.
\]
取 $n=2018,x=\pi/6$，则 $(n+1/2)x=4037\pi/12$ 与 $5\pi/12$ 相差 $168\cdot2\pi$。令 $\theta=\pi/12$，有 $\sin(5\pi/12)=\cos\theta$、$\cos(5\pi/12)=\sin\theta$，因此
\[
\sum_{j=1}^{2018}j\sin(j\pi/6)
=-\frac{4037}{4}+\frac14\cot^2\theta
=-\frac{4037}{4}+\frac{(2+\sqrt3)^2}{4}
=\boxed{\sqrt3-\frac{2015}{2}}.
\]
\end{proof}
\textbf{另一种核对。}每连续 $12$ 项正弦值之和为零，而 $\sum_{j=1}^{12}j\sin(j\pi/6)=-12-6\sqrt3$。故前 $2016=168\cdot12$ 项的加权和为 $168(-12-6\sqrt3)$；加上末两项 $2017/2+1009\sqrt3$，也得到相同结果。权 $j$ 本身没有周期，但平移一个完整周期所多出的常数权乘上零和，恰好抵消。

\section{方法之间的联系}\label{sec:methods}
\subsection{计数：选择能够唯一恢复对象的中间量}
2009 年高考固定二次方程中的一个系数，2015 年南京、盐城一模固定子集最大元，2018 年高考固定最大元素的插入位置，表面不同，核心都在于建立可逆分类。分类标准必须使每个原对象恰落入一类，类内各个选择再由乘法原理计算。检查公式时先算最小规模，是发现漏项和重计的有效方式，但有限个数值吻合不能代替一般证明。

另一些题通过双射比较两类对象，而不必分别算出它们。2012 年高考中，每条倍增链只需决定首项颜色；2014 年南京、盐城一模中，删去零组后只需决定第一个非零组的符号；2016 年三市二模中，全部位取反把余数 $1$ 与余数 $3$ 对应。解题时应写明变换的逆操作或唯一恢复过程，从而说明“没有漏、没有重”。

\subsection{概率：先确定样本空间，再决定是否压缩状态}
2017 年高考抽球题的适当样本空间是黑球占据的位置集合；2019 年高考格点题的样本空间是无序点对。二者都采用等可能计数，却不能把概率分母随意换成有序数对或独立试验总数。独立性只在题目允许时使用，不能因为每一步边缘概率相同就断言各步独立。

图上游走则适合把多个顶点合并为少数状态。合并的依据是同一状态内的任一顶点转移到另一状态的概率都相同。2011 年四面体题中，三个非出发点到出发点的概率均为 $1/3$；2013 年三棱柱题中，同一底面的三个顶点到另一底面的概率也均为 $1/3$。若这些转移概率不同，仅记录“在哪一类”就可能不足以推出递推，必须细分状态。

\subsection{代数与几何：寻找运算后保持形式的表示}
2010 年高考用余弦递推保持有理数运算，2014 年高考对 $xf_0(x)=\sin x$ 求导保持相邻导数组合，2015 年泰州二模通过 $x\leftrightarrow1/x$ 保持多项式形式。共同策略是先观察所求对象的结构，再选恒等式；不必总先求出每个中间量的显式通项。

抛物线题中，参数表示 $(t^2,2t)$ 同时简化曲线方程、切线方程和弦方程。对于 $y^2=4x$，两切点参数的和与积分别控制切线交点的纵、横坐标，而弦方程为 $2x-(s+t)y+2st=0$。第~\ref{sec:c2010nj1}、\ref{sec:c2012nty1}、\ref{sec:c2014sb} 节分别利用乘积固定、交点坐标相同、弦斜率固定。使用这一表示时仍须保留两点不同、参数非零等原题条件，不能因式子简洁而忽略退化情形。

\section{题目完整但市级归属待核的材料}\label{sec:unattributed}
本节材料不能计入任何城市的一模、二模或三模完成数。保留独立题解，是为了避免数学内容与来源核验互相混淆；只有补到足以确认命题单位的原卷首页或同期发布记录，才能迁入市级正文。
\subsection{2019 届“第二次模拟考试”材料：同色三元组的最小数}\label{sec:u2019}
\textbf{来源与限制。}\href{https://dfiles.jiajiaoban.com/ohr/2020/03/23/120155_5e7834b332ec4.pdf}{公开 PDF：2019 届高三年级第二次模拟考试}，附加题第 23 题。取得的文本含完整题干，但首页未标明城市；文件上传时间也不能代替考试年份或命题单位，暂不认定为南京、盐城或苏锡常镇卷。
\begin{example}[题意]
平面上 $2n$ 个不同点分别染为红色或蓝色，$n\geq3$。从中任取三个点，三个点颜色相同的取法总数记为 $T$。求 $n=3$ 时 $T$ 的最小值；当 $n\geq4$ 时，证明 $T\geq2\binom n3$。
\end{example}
\textbf{思路。}本题数的是三点子集，不要求三点构成三角形，点的位置不影响答案。用红点数 $r$ 表示全部计数，再把偏离均分的程度写成一个平方。
\begin{proof}
若红点有 $r$ 个，蓝点有 $2n-r$ 个，则
\[
T=\binom r3+\binom{2n-r}3,
\qquad 0\leq r\leq2n.
\]
利用 $\binom s3=s(s-1)(s-2)/6$（对整数 $s=0,1,2$ 也成立），展开并整理得
\[
T=2\binom n3+(n-1)(r-n)^2.
\]
因为 $n-1>0$，最小值为 $2\binom n3$，当且仅当两色各有 $n$ 个点时达到。特别地，$n=3$ 时最小值为 $2$；对 $n\geq4$，所需不等式也立即成立。
\end{proof}
\textbf{复盘。}“二染色三点集合”容易让人误想到边染色的拉姆齐问题。本题染的是点，颜色相同的三点数只取决于两种颜色各有多少个点，不涉及边色，也不需要假设任意三点不共线。

\section{城市索引与检索登记}\label{sec:register}
\subsection{本稿实际覆盖情况}
本稿高考部分列出 2008—2019 年的 12 道附加题末题，其中 2008 年最后一式、2014 年第二问保留了原文符号待核的校勘说明。市级正文目前有 15 道不重复的末题；联考分别映射到参与城市后，涉及全部 13 市，但这\keyterm{不等于 13 市各年各轮已收齐}。另有 1 道 2019 届材料已解、城市归属待核，不计入市级题解数。本次检索登记日期为 2026 年 10 月 11 日。

下表只索引已写入市级正文的题目。除这些具体条目外，2008—2019 年每市其余一模、二模、三模槽位均未完成收录；后表列出部分已经找到的进一步线索。城市列表不是对联考组合长期不变的假设：例如某年三市同卷，不能推出另一年仍是这三市同卷。
\begin{longtable}{p{0.13\textwidth}p{0.78\textwidth}}
\textbf{城市}&\textbf{已经写入正文的试卷及题解位置}\\\hline
南京&2010 一模（第~\ref{sec:c2010nj1} 节）；2011 三模（第~\ref{sec:c2011ny3} 节）；2013 三模（第~\ref{sec:c2013ny3} 节）；2014 一模（第~\ref{sec:c2014ny1} 节）；2015 一模（第~\ref{sec:c2015ny1} 节）。\\
无锡&2013 苏锡常镇一调（第~\ref{sec:c2013sxcz1} 节）。\\
徐州&2014 届苏北四市期末调研（第~\ref{sec:c2014sb} 节）。\\
常州&2013 苏锡常镇一调（第~\ref{sec:c2013sxcz1} 节）。\\
苏州&2013 苏锡常镇一调（第~\ref{sec:c2013sxcz1} 节）。\\
南通&2010 二模（第~\ref{sec:c2010nty2} 节）；2012 一调（第~\ref{sec:c2012nty1} 节）；2016 二模（第~\ref{sec:c2016nty2} 节）；2017 一模（第~\ref{sec:c2017nt1} 节）；2018 一模（第~\ref{sec:c2018nttz1} 节）。\\
连云港&2014 届苏北四市期末调研（第~\ref{sec:c2014sb} 节）。\\
淮安&2014 届苏北四市期末调研（第~\ref{sec:c2014sb} 节）。\\
盐城&2011 三模（第~\ref{sec:c2011ny3} 节）；2013 三模（第~\ref{sec:c2013ny3} 节）；2014 一模（第~\ref{sec:c2014ny1} 节）；2015 一模（第~\ref{sec:c2015ny1} 节）；2018 三模（第~\ref{sec:c2018yc3} 节）。\\
扬州&2010 二模（第~\ref{sec:c2010nty2} 节）；2012 一调（第~\ref{sec:c2012nty1} 节）；2016 二模（第~\ref{sec:c2016nty2} 节）。\\
镇江&2013 苏锡常镇一调（第~\ref{sec:c2013sxcz1} 节）；2018 镇江一模（第~\ref{sec:c2018zj1} 节）。\\
泰州&2010 二模（第~\ref{sec:c2010nty2} 节）；2012 一调（第~\ref{sec:c2012nty1} 节）；2015 泰州二模（第~\ref{sec:c2015tz2} 节）；2016 二模（第~\ref{sec:c2016nty2} 节）；2018 一模（第~\ref{sec:c2018nttz1} 节）。\\
宿迁&2014 届苏北四市期末调研（第~\ref{sec:c2014sb} 节）。\\\hline
\end{longtable}

\subsection{按年检查缺口}
下表统计独立题目，不按联考城市重复计数。“0”仅表示本稿尚未完成该年市级题解，绝不表示该年没有模拟考试。
\begin{longtable}{p{0.1\textwidth}p{0.13\textwidth}p{0.66\textwidth}}
\textbf{年份}&\textbf{市级题数}&\textbf{已收录范围与主要缺口}\\\hline
2008&0&起始年份的市级卷尚缺；已找到南京二模同期报道，尚未取得末题可靠文本。\\
2009&0&苏北四市三调取得部分答案，题干不完整；苏锡常镇一调只有发布线索。\\
2010&2&南京一模、通扬泰二模已收；其余轮次及苏锡常镇、苏北卷仍缺。\\
2011&1&南京、盐城三模已收；南京一模公式、苏锡常镇一调和二调末题仍待核。\\
2012&1&通泰扬一调已收；苏锡常镇二调条件符号待核，其他轮次仍缺。\\
2013&2&苏锡常镇一调、南京盐城三模已收；宁盐一模及苏中二模公式仍待核。\\
2014&2&宁盐一模、苏北期末调研已收；苏锡常镇二模只取得末题答案，需补题干。\\
2015&2&宁盐一模、泰州二模已收；苏锡常镇两轮、盐城三模等仍缺末题可靠文本。\\
2016&1&通扬泰二模已收；苏锡常镇二模公式缺损，其他地区和轮次仍缺。\\
2017&1&南通一模已收；宁盐一、二模，六市二模、苏锡常镇一模等仍待补。\\
2018&3&镇江一模、通泰一模、盐城三模已收；宁盐联考、六市二模、苏锡常镇二模等仍待补。\\
2019&0&取得多套市级卷发布页，但尚未核清其末题；另列归属待核题 1 道。\\\hline
\end{longtable}

\subsection{可继续追溯的待核材料}
以下均为\keyterm{线索登记}，不是已收录题目目录。题名以检索到的页面为线索，可能仍有错标；每条说明目前缺少什么。对只找到参考答案的材料，不按答案倒造一份自称“原题”的题干；对转录中消失的分式、上下标和不等号，也不靠猜测补全后混入正式题库。
\begin{enumerate}
\item \href{https://edu.sina.com.cn/gaokao/2008-04-28/1038136658.shtml}{2008 南京二模同期报道}：确认当年数学附加卷的考试背景，尚非末题原文。
\item \href{https://doc.xuehai.net/b8dc14801260df2bee36b2ffd.html}{2009 苏北四市第三次调研}：首页列徐州、宿迁、淮安、连云港及 2009 年 3 月 31 日；后页有排列数求和、递推与连乘积答案，末题完整题干未核得。保留“三调”原名，不因名称中的“三”擅自改成通常意义的五月三模。
\item \href{https://m.liuxue86.com/a/485326.html}{2009 苏锡常镇教学情况调查（一）}：已找到试卷发布条目，末题正文尚缺。
\item \href{https://www.tthaoke.com/shijuan_37108.html}{2010 苏锡常镇教学情况调查（一）}：目录单列附加卷及答案 PDF；本次未取得可核验末题。网页上传年份与文件试卷年份不同，按文件题名登记为 2010。
\item \href{https://doc.xuehai.net/ba99b4701970c23149844d317.html}{题名为 2010 南京二模的转载}：末题与本稿南京一模抛物线题相同，存在轮次归属冲突；不重复计数，需核原卷首页。
\item \href{https://doc.xuehai.net/bdb0bd5ffdb8e57341a9c50cf.html}{2010 苏锡常镇三模线索}：可读部分与图片混排，尚未核得附加题末题；不能把主卷函数压轴题移入本书。
\item \href{https://doc.xuehai.net/bd5dae6b568cf4072bb664ad4.html}{2011 南京一模}：数列末题中的商与不等号在转录中缺损，需要核对关系符号及归纳证明要求。
\item \href{https://wenku.baidu.com/view/47100f08581b6bd97f19eaf9.html}{2011 苏锡常镇第一次调研}及\href{https://max.book118.com/html/2016/1021/60264575.shtm}{另一转载}：可定位主卷，尚未取得完整附加题末题。
\item \href{https://doc.xuehai.net/b5e02fd9333058500970cd509.html}{2011 苏锡常镇二模答案}：图片型材料，末题未完成校读。
\item \href{https://max.book118.com/html/2017/0331/98187113.shtm}{2011 通扬泰二模线索}：只有题名及局部预览，未计入正文。
\item \href{https://www.renrendoc.com/paper/112511852.html}{2011 扬州三模相关页面}：取得的内容不足以确认附加题末题的完整条件。
\item \href{https://www.1mpi.com/doc/9074c29f2b110b8bb634e626}{2012 苏锡常镇一调}：公开文本末尾仍为主卷“曲线宽度”题，没有取得附加题，不能以该末尾题替代。
\item \href{https://www.mianfeiwendang.com/doc/fbfc7ec96dd36d63a7bf9e7d/5}{2012 苏锡常镇二调}：已定位第 23 题，涉及 $f_n(x,y)=(x+y)^n-(x^n+y^n)$ 与数学归纳法；参数 $a,b$ 的分式约束和目标不等号未完整保留，暂缓收录。
\item \href{https://doc.xuehai.net/b82f390d6798cd90b0ed7c51f.html}{2012 通扬泰二模}：末题对数函数及乘积不等式中的正负号缺损，不能可靠恢复题目。
\item \href{https://www.51wendang.com/doc/517100508630501f27d3bf9d}{2012 南京、盐城二模线索}：公开部分不足以核清末题。
\item \href{https://doc.xuehai.net/b5cafa3f91746899ae088795f.html}{2012 南京三模线索}：本次未能读取完整文档，需另找原卷或同期转载。
\item \href{https://www.51wendang.com/doc/30287572122302d8cbc36f9e/7}{2013 南京、盐城一模}：第 23 题涉及二项式展开和含根式数的幂，但关键公式缺损，答案后页读取失败。
\item \href{https://doc.xuehai.net/b9a2cc392ff2ad5ca2fdcb82f.html}{2013 苏中四市二模}：取得涉及 $x+1/x$ 与幂比较的部分文字；原函数定义和全部小问尚不完整，联考具体组成也需复核。
\item \href{https://max.book118.com/html/2017/0331/98180636.shtm}{2014 苏锡常镇一调线索}：预览未核得末题全部条件。
\item \href{https://doc.xuehai.net/bd211a532a3be38c0302b5df3-8.html}{2014 苏锡常镇二模答案}：第 23 题答案出现 $a_2=6,a_3=6,a_4=18$ 和相邻点标数递推，符合环上三种标号的模型；但只有答案，尚缺原题完整定义，不据此倒造原题。
\item \href{https://doc.xuehai.net/b11a07d49cb418f9a37ccf351.html}{2014 盐城三模线索}：可见部分不足以核得附加题末题。
\item \href{https://doc.xuehai.net/babe53a180f948f2c555dad51.html}{2015 苏锡常镇一调}：后续页以图片为主，尚缺可核验的末题文本。
\item \href{https://doc.xuehai.net/b9341d44caa5719f543ac937a.html}{2015 苏锡常镇二调}及\href{https://www.mianfeiwendang.com/doc/8ec448b5e1bc8af6e7874d05}{另一转载}：后者止于主卷第 20 题，前者附加题信息亦未核清。
\item \href{https://www.51wendang.com/doc/b354a966ae74314766a3ac38}{2014 宁盐一模页面的相关文档目录}列有“盐城市 2015 届第三次模拟考试数学试题及答案”：关联文档可读部分只到主卷，尚缺附加题。此项只是目录线索。
\item \href{https://max.book118.com/html/2016/0412/40228345.shtm}{2016 苏北四市期末调研}：公式转录缺损，末题暂不能可靠重述。
\item \href{https://www.renrendoc.com/paper/195156858.html}{苏锡常镇教学情况调研（二）混合转载}：正文含 2015--2016 试卷，同时夹有 2013--2014 答案标题，存在混卷风险；不按网页标题直接认定题目年份。
\item \href{https://gaokao.eol.cn/shiti/moniti/201603/t20160325_1380484.shtml}{2016 南京、盐城二模发布页}：已定位发布信息，末题仍待取得。
\item \href{https://imgs.app.gaokaozhitongche.com/resource/pdf/paper/202305/20230504165103_LpudIQ.pdf}{2017 南京、盐城一模 PDF 线索}：本次文字提取未取得可用公式，不声称已经完成原卷校读。
\item \href{https://app.gaokaozhitongche.com/news/10648}{2017 南京、盐城二模发布页}：公开内容不足以核清末题。
\item \href{https://gaokao.eol.cn/shiti/sx/201705/t20170510_1514413.shtml}{2017 南京三模发布页}：尚缺附加题末题的可靠转录。
\item \href{https://www.renjiaoshe.com/gaozhongshuxue/yilun/172495.html}{2017 苏锡常镇一模线索}：网页读取报错，未取得题干。
\item \href{https://news.koolearn.com/20170324/1115209.html}{2017 苏北六市二模汇总}及\href{https://news.koolearn.com/20170323/1115099.html}{数学分卷入口}：六市为南通、扬州、泰州、徐州、淮安、宿迁；数学分卷本次未能完整读取。
\item \href{https://imgs.app.gaokaozhitongche.com/resource/pdf/paper/202305/20230505003302_0So7fy.pdf}{2018 南京、盐城一模 PDF 线索}：未提取到可核验的末题文本。
\item \href{https://news.koolearn.com/20180321/1146694_2.html}{2018 南京、盐城二模发布页}：尚缺附加题末题的可靠转录。
\item \href{https://www.renjiaoshe.com/gaozhongshuxue/erlun/508681.html}{2018 六市第二次调研}：页面索引明确含附加题，但正文读取报错；不把同一地区 2016 年三市卷当作本卷。
\item \href{https://news.koolearn.com/20180320/1146651.html}{2018 苏锡常镇一模汇总}及\href{https://news.koolearn.com/20180503/1149614.html}{二模汇总}：目前仅登记发布线索。另有\href{https://max.book118.com/html/2019/0914/7103142025002055.shtm}{二模文档转载}，末题尚待核验。
\item \href{https://m.koolearn.com/news/20190115/1167775.html}{2019 南京、盐城一模发布页}及\href{https://zj.qinxue100.com/news/gz/112567.html}{参考答案页}：已定位同期发布记录，末题公式尚未核全。
\item \href{https://m.koolearn.com/gaokao/20190321/1211493.html}{2019 南京、盐城二模发布页}：另有\href{https://max.book118.com/html/2019/0419/8102132114002017.shtm}{标为“五套精选”的转载}，该转载不能不加甄别地整体视为一份市级卷。
\item \href{https://app.gaokaozhitongche.com/news/h/0O7gr6OA}{2019 南京三模发布页}：可确认附加题材料入口，公开内容未完成题干校读。
\item \href{https://m.koolearn.com/news/20190319/1170460.html}{2019 苏锡常镇一模线索}及\href{https://www.renrendoc.com/p-83313379.html}{二模转载}：后者末题含函数列 $f_n,g_n$，定义公式缺损，不据残文补题。
\item \href{https://max.book118.com/html/2019/0419/7121015136002021.shtm}{2019 苏北四市一模线索}及\href{https://max.book118.com/html/2023/1006/7140025105005164.shtm}{另一个苏北三市一检条目}：两者题名和联考组成不同，不能未经核对就合并；末题均未完成核验。
\item \href{https://max.book118.com/html/2019/0630/8101117122002032.shtm}{2019 七市三模线索}：参与城市、附加题和参考答案尚需原卷确认，不外推到所有江苏城市。
\end{enumerate}

\section{制度与高考原题来源}\label{sec:sources}
以下优先列出考试院或同期发布材料；使用转载时明确注明。来源提供题目和历史事实的依据，本书解答为重新组织的推导。网页中的文字提取可能丢失数学符号，具体校勘说明已随题列出；不能把“页面有标题”与“原题每个符号已核清”混为一谈。
\begin{enumerate}
\item\label{src:history} \href{https://gaokao.chsi.com.cn/gkxx/ss/200604/20060405/432170.html}{阳光高考：江苏自主命题相关报道（2006 年）}，用于核对江苏自 2004 年起自主命题的历史起点。
\item\label{src:scheme} \href{https://www.jseea.cn/webfile/highschool_ty_artxk_files/2007-05-24/3187.html}{江苏省教育考试院：江苏省 2008 年普通高考方案}；\href{https://gaokao.chsi.com.cn/gkxx/js/200711/20071121/1598706.html}{阳光高考：2008 年高考考试说明出炉}；\href{https://gaokao.chsi.com.cn/gkxx/js/200711/20071121/1598761.html}{考试说明公布消息}。用于区分自主命题起点与 $160+40$ 试卷结构起点。
\item\label{src:g2008} \href{https://www.yeyulingfeng.com/wendang/194472.html}{2008 年江苏高考数学解析卷转载}，末题在转载中重排为第 26 题；另参\href{https://doc.xuehai.net/b69540a89a069c31e3680dd29.html}{阙东进《2008 年江苏高考第 23 题杂谈》转载}，《中学数学月刊》2009 年第 3 期，第 45 页起。最后一个积分恒等式的转录有疑点，正文同时给出两种求和下限对应的正确值，尚不能视为原卷逐字定稿。
\item\label{src:g2009} \href{https://www.jseea.cn/webfile/gaokao/2009-06-09/6597.html}{江苏省教育考试院：2009 年试题发布入口}；\href{https://29805199.s21i.faiusr.com/61/3/ABUIABA9GAAgqP6QowYovoST9wI.pdf}{2009 年数学试题及答案 PDF 转载}；\href{https://www.yeyulingfeng.com/wendang/15424.html}{另一文本转载}。
\item\label{src:g2010} \href{https://gaokao.chsi.com.cn/gkxx/ss/201006/20100612/96768571.html}{阳光高考：2010 江苏数学试题相关发布}；\href{https://www.yeyulingfeng.com/wendang/16407.html}{2010 年江苏数学试卷转载}。原卷第 23 题为余弦有理性问题，不能混同转载重排的概率题号。
\item\label{src:g2011} \href{https://www.yeyulingfeng.com/wendang/202979.html}{2011 年江苏高考数学试卷转载}，整数点差与同余计数题。
\item\label{src:g2012} \href{https://www.yeyulingfeng.com/wendang/21104.html}{2012 年江苏高考数学试卷转载}，集合中元素与其二倍的互斥归属题。
\item\label{src:g2013} \href{https://www.jseea.cn/webfile/upload/old/contents/channel_6/2013/06/1306100857696.pdf}{江苏省教育考试院：2013 年高考试题及参考答案汇编 PDF}，数学附加题第 23 题。
\item\label{src:g2014} \href{https://mat1.gtimg.com/edu/PDF/jsws.pdf}{2014 年江苏高考数学试题及答案 PDF 转载}；\href{https://www.renrendoc.com/paper/207419020.html}{含附加题的详细答案文本转载}；\href{https://gaokao.eol.cn/shiti/zhenti/201406/t20140607_1127957.shtml}{中国教育在线同期发布入口}。第二问的绝对值符号仍按正文说明保留校勘状态。
\item\label{src:g2015} \href{https://www.jseea.cn/webfile/upload/old/contents/channel_6/2015/06/1506101049510.pdf}{江苏省教育考试院：2015 年高考试题及参考答案汇编 PDF}，数学附加题第 23 题。
\item\label{src:g2016} \href{https://www.jseea.cn/webfile/upload/old/contents/channel_173/2016/06/1606100927946.pdf}{江苏省教育考试院：2016 年高考试题及参考答案汇编 PDF}，数学附加题第 23 题。
\item\label{src:g2017} \href{https://www.jseea.cn/webfile/upload/old/contents/channel_1/2017/06/1706092140912.pdf}{江苏省教育考试院：2017 年高考试题及参考答案汇编 PDF}，数学附加题第 23 题。
\item\label{src:g2018} \href{https://www.jseea.cn/webfile/upload/old/contents/channel_173/2018/06/1806100757613.pdf}{江苏省教育考试院：2018 年高考试题及参考答案汇编 PDF}，数学附加题第 23 题。
\item\label{src:g2019} \href{https://www.jseea.cn/webfile/highschool_ptgk_data/2019-06-10/6514.html}{江苏省教育考试院：2019 年高考试题发布入口}；\href{https://gaokaohub.com/q/item/102459044642881.html}{2019 年格点距离末题的公式文本转载}。
\end{enumerate}
\end{document}
